\documentclass[CMPE]{KGCOEReport}

\usepackage{graphicx}
\usepackage{newtxtext, newtxmath}
\usepackage{parskip}
\usepackage{karnaugh-map}
\usepackage{float}
\usepackage{pdfpages}

\graphicspath{{./assets/}}

\newcommand{\classCode}{CMPE 250}
\newcommand{\name}{Aden Perry}
\newcommand{\LabSectionNum}{L1}
\newcommand{\LabInstructor}{Chongzhou Fang}

\newcommand{\TAs}{Eddie Satowski \\ Eric Reed \\ Khoi Pham}
\newcommand{\LectureSectionNum}{01}
\newcommand{\LectureInstructor}{Chongzhou Fang}
\newcommand{\exerciseNumber}{7}
\newcommand{\exerciseDescription}{Circular FIFO Queue Operations}
\newcommand{\dateDone}{2026-03-05}
\newcommand{\dateSubmitted}{2025-03-19}

\begin{document}

\maketitle

\section*{Abstract}

The purpose of this exercise was to desgin a circular queue using the NXP
Freedom KL05Z board. The queue was described with a record in memory, and
implemented inside of a buffer in memory. Subroutines took the record address
as arguments and performed I/O on the buffer. A TUI menu was created to control
the queue; string and character I/O subroutines were reused from previous
exercises. The TUI allowed users to enqueue, dequeue, print queue, print queue
status, and get help. The program was translated, compiled, and flashed to the
KL05Z. It was tested with various inputs at various queue capacities. All
outputs were printed correctly. The queue properly handled dequeue from empty,
enqueue to full, and displaying queue contents in decimal and hexadecimal
representations. The exercise was successful.

\section*{Procedure}

A Keil MDK-ARM project was created. Provided libraries were added to Source
Group 1 of the project. A subroutine InitQueue was written to create a record
structure describing the queue. The structure contained six fields: two 4-byte
data pointers, two 4-byte buffer pointers, a 1-byte buffer size, and a 1-byte
queue size. The data pointers were used to access queue data. The buffer
pointers were used to bound the edges of the buffer. InitQueue took in a buffer
start address, a buffer capacity in bytes, and a record address.

A subroutine Dequeue was written to get data from the queue. It took in the
record address as input and gave a dequeued byte as output. If a dequeue
operation was attempted on an empty queue, the PSR C flag was set. The
subroutine compared the queue size to zreo to determine if the queue was empty.
It read the data at the address specified in the data out pointer of the record
to register R0. It incremented the data out pointer after use; if the data out
pointer exceeded buffer bounds, it was set to the start of the buffer. The
subroutine decremented the queue size contained in the record.

A subroutine Enqueue was written to put data into the queue. It took in the
record address and the byte to enqueue as inputs. If an enqueue operation was
attempted on a full queue, the PSR C flag was set. The subroutine compared
buffer size and queue size to determine if the queue was full. It stored byte
0 of register R0 to the address contained the data in pointer of the record. It
incremented the data in pointer after use. The data in pointer was incremented
in the same manner as the data out pointer. The subroutine incremented the
queue size contained in the record.

A subroutine PutNumUB was written as a wrapper around the PutNumU subroutine
implemented in a previous exercise. PutNumUB extracted the least significant
byte of an input and passed the byte as an argument to PutNumU, which converted
the byte to decimal and sent it to the Universal Asynchronous
Receiver/Transmitter (UART) in ASCII encoding.

A subroutine PutNumHex was written to convert a register to hexadecimal and
send it in ASCII encoding to the UART. It took in a register as input. The
subroutine iteratively masked and converted each nibble to a hexadecimal ASCII
value. The mask position was adjusted by with logical right shifting. The ASCII
value was sent to the PutChar subroutine written in a previous exercise.

A TUI was implemented to request input from a user. The user was presented with
five options: dequeue, enqueue, help, print, and status. The dequeue option
attempted to dequeue an item using the Dequeue subroutine. If successful, it
sent the dequeued character to the UART. The enqueue option attempted to
enqueue an item using the Enqueue subroutine. The user was asked for a
character to enqueue. The help option printed a string listing available
options. The print option option printed the entire queue without modifying it.
The status option printed the information contained in the record. Status
information was also printed after enqueue and dequeue operations.

\section*{Results}

The TUI was tested using a variety of inputs. The complete test is given by
Figure \ref{fig:demo}. It demonstrates each command under a variety of
conditions. During the test, non-command characters were entered, but were
successfully ignored.

\begin{figure}[H]
	\center
	\includegraphics[height=0.95\textheight]{assets/sc}
	\caption{Demonstration of TUI}
	\label{fig:demo}
\end{figure}

As seen in Figure \ref{fig:demo}, commands enqueue and dequeue were tested
during queue full, queue empty, and queue partially full states. Each time, the
command printed correctly the result of the operation, as well as the status
of the queue record. Also shown is that the queue stays contained within the
boundaries of the queue buffer: the data in and out pointers jump from
\texttt{0x1FFFFD03} to \texttt{0x1FFFFD00}, because they have hit the edge of
the buffer.

A useful way by which to analyze a given assembly (or other) program is by
examining its memory map. The memory map for the implemented program is
presented in Table \ref{tab:mem}.

\begin{table}[H]
	\center
	\caption{Memory Ranges}
	\begin{tabular}{|c|c|c|}
		\hline
		Symbol Name  & Start Address & End Address \\
		\hline
		\hline
		MyCode       & 0x00000410    & 0x000007a7  \\
		\hline
		CMD\_S       & 0x000001cc    & 0x000001ed  \\
		\hline
		CRLF\_S      & 0x000001ee    & 0x000001f0  \\
		\hline
		FAIL\_S      & 0x000001f1    & 0x00000201  \\
		\hline
		CHAR\_S      & 0x00000202    & 0x00000211  \\
		\hline
		STAT\_S      & 0x00000212    & 0x00000222  \\
		\hline
		SUCC\_S      & 0x00000223    & 0x00000233  \\
		\hline
		ENQU\_S      & 0x00000234    & 0x00000244  \\
		\hline
		HELP\_S      & 0x00000245    & 0x00000280  \\
		\hline
		STAT\_IN\_S  & 0x00000281    & 0x0000028a  \\
		\hline
		STAT\_OUT\_S & 0x0000028b    & 0x00000295  \\
		\hline
		STAT\_NUM\_S & 0x00000296    & 0x0000029f  \\
		\hline
		QBUFFER      & 0x1ffffd00    & 0x1ffffd4f  \\
		\hline
		QRECORD      & 0x1ffffd50    & 0x1ffffd61  \\
		\hline
	\end{tabular}
	\label{tab:mem}
\end{table}

The MyCode symbol represents all executable code inside the program. According
to Table \ref{tab:mem}, the MyCode area occupies 920 bytes of instructions.
Thiss includes the main TUI program, as well as the command selection and
execution subroutines. Various string constants were created to enable string
printing. In total, they occupy 212 bytes of memory. Additionally, two
variables were created in RAM to store the queue and the queue record. The
queue record occupies 18 bytes of space. The queue buffer occupies 80 bytes of
space. Thus, although only 4 bytes were used in program testing, this program
could support a queue of up to 80 bytes in length without remapping RAM.

\section*{Conclusion}

The objective of this exercise was to design a circular FIFO queue, as well as
to design  subroutines that could print a byte in decimal and a register in
hexadecimal to the UART. The queue used a record to manage meta information,
and used a buffer to store queue contents. In this program, the queue was used
to facilitate a TUI menu, but in future exercises it may have more practical
applications. For example, a queue may be used for interrupt transmition with
the UART<F10>. Values to transmit may be enqueued, to be transmitted at the
leisure of the connected device.

\includepdf[pages=-]{assets/prelab}

\end{document}
